\documentclass[twocolumn]{article}
\usepackage[T1]{fontenc}
\usepackage{amsmath}
\usepackage{graphicx}

\title{Capacitor Charging}
\author{Clyde Fernandez}
\date{}

\begin{document}
\maketitle

\begin{abstract}
This document introduces the charging behaviour of a resistor-capacitor (RC) network and derives the capacitor voltage at one time constant.    
\end{abstract}

\section{Introduction}
This document introduces the charging behavior of a resistor-capacitor (RC) network. When a capacitor C is charged through a series R from a constant supply $V_0$, the voltage of the capacitor increases according to the condition.
\begin{equation}
\label{one}
    V(t) = V_0 \left(1-e^{-t/RC}\right),
\end{equation}
where the product RC has units of seconds and is called the time constant $\tau = RC$.

\section{Voltage at the Time Constant}
Substituting $t = \tau$ into (1):

\begin{align}
V(\tau) &= V_0\left(1-e^{-\tau/RC}\right) \notag \\
        &= V_0\left(1-e^{-RC/RC}\right) && \text{since } \tau = RC \notag \\
        &= V_0\left(1-e^{-1}\right) \label{eq:vtau} \\
        &\approx V_0(1-0.368) \notag \\
        &\approx 0.632\,V_0. \notag
\end{align}

\noindent The exponent cancelling to $-1$ in (2) is why the $0.632$ fraction is
universal --     $R$ and $C$ dropped out of the calculation before the final
line, so this fraction holds for any $RC$ network. Substituting
$V_0 = 5.0\,\mathrm{V}$:
\begin{equation}
\label{three}
  V(\tau) \approx 0.632 \times 5.0\,\mathrm{V} \approx 3.16\,\mathrm{V},
\end{equation} 
matching the dashed line in Figure 1.

\begin{figure}[t]
\centering
\includegraphics[width=\columnwidth]
{rc_charging.pdf}
\caption{Charging curve for
$R=10\,\text{k}\Omega$, $C=22\,\text{pF}$.}
\label{fig:rc}
\end{figure}



\end{document}
